\documentclass[10pt,letterpaper]{article}
\usepackage{cornell-notes}

\title{Chapter 4: File Systems}
\author{}
\date{}

\setDocTitle{Chapter 4: File Systems}
\setDocSubtitle{Operating Systems}
\setDocVersion{v1.0}
\setDocStatus{Study Companion}
\setDocOwner{Operating Systems Cornell Notes}
\setDocAuthor{}
\setDocDate{\today}
\setDocSummary{Cornell-format study and review notes for Chapter 4: File Systems.}

\setCornellCollection{Operating Systems Cornell Notes}
\setCornellUnitType{Chapter}
\setCornellUnitNumber{4}
\setCornellUnitTitle{File Systems}
\setCornellCompanionLabel{Study and review companion}

\hypersetup{pdftitle={Chapter 4: File Systems},pdfsubject={Cornell notes},pdfauthor={}}

\begin{document}
\maketitle
\section*{Learning Objectives}
\begin{studybox}{By the end of this chapter, you should be able to}
\begin{enumerate}
  \item Explain file naming, structure, type, access, attributes, and operations.
  \item Resolve absolute and relative paths through hierarchical directories.
  \item Compare contiguous, linked, table-based, and indexed file allocation.
  \item Explain directory implementation, links, virtual file systems, journaling, and log-structured designs.
  \item Evaluate free-space management, backup, consistency checking, caching, read-ahead, and placement.
  \item Compare the representative allocation and directory choices in MS-DOS, UNIX V7, and CD-ROM file systems.
\end{enumerate}
\end{studybox}

\begin{studybox}{Section Map}
\hyperlink{sec-4-1}{4.1} \quad \hyperlink{sec-4-2}{4.2} \quad \hyperlink{sec-4-3}{4.3} \quad \hyperlink{sec-4-4}{4.4} \quad \hyperlink{sec-4-5}{4.5} \quad \hyperlink{sec-4-6}{4.6} \quad \hyperlink{sec-4-7}{4.7}
\end{studybox}

\section*{Cornell Cue-and-Notes Area}
\small
\begin{longtable}{C N}
\rowcolor{navy}
\color{white}\textbf{Cue / Recall Prompt} &
\color{white}\textbf{Notes}\\
\endfirsthead
\rowcolor{navy}
\color{white}\textbf{Cue / Recall Prompt} &
\color{white}\textbf{Notes (continued)}\\
\endhead
\rowcolor{ice}\multicolumn{2}{l}{\hypertarget{sec-4-1}{}\textbf{Section 4.1: Files}}\\*\hline
\textbf{File abstraction} & A file is a named persistent information container. The interface separates logical content and metadata from sectors, controllers, and physical placement.\\\hline
\textbf{4.1.1 Naming} & Names may be case sensitive, length limited, or constrained by syntax. Extensions can guide applications but need not determine the file's actual content.\\\hline
\textbf{4.1.2 Structure} & Systems may treat a file as an unstructured byte sequence, a sequence of records, or a tree of records. A byte stream gives applications the greatest freedom.\\\hline
\textbf{4.1.3 Types} & Regular files, directories, character and block special files, and other special objects support different semantics. Magic numbers or metadata can identify format.\\\hline
\textbf{4.1.4 Access} & Sequential access consumes data in order; random access selects an offset or record directly. The file interface tracks a current position when appropriate.\\\hline
\textbf{4.1.5 Attributes} & Metadata can include size, owner, protection, timestamps, flags, and storage information. The kernel maintains attributes separately from file content.\\\hline
\textbf{4.1.6 Operations} & Create, delete, open, close, read, write, append, seek, inspect, rename, and change attributes form the core file-operation set.\\\hline
\textbf{4.1.7 Call sequence} & A typical program opens a path, checks the returned descriptor, loops over read or write results, handles partial transfers and errors, and closes the descriptor.\\\hline
\rowcolor{ice}\multicolumn{2}{l}{\hypertarget{sec-4-2}{}\textbf{Section 4.2: Directories}}\\*\hline
\textbf{4.2.1 Single-level} & One directory provides simple naming but causes collisions and weak grouping as the number of users and files grows.\\\hline
\textbf{4.2.2 Hierarchies} & A rooted tree allows directories to contain files and subdirectories. Per-user subtrees and mount points create scalable namespaces.\\\hline
\textbf{4.2.3 Path names} & Absolute paths begin at the root; relative paths begin at the working directory. Components such as current and parent directories affect traversal.\\\hline
\textbf{4.2.4 Operations} & Create and remove directories, open and read entries, rename, link, unlink, and change the working directory. Removal policies must address nonempty directories.\\\hline
\rowcolor{ice}\multicolumn{2}{l}{\hypertarget{sec-4-3}{}\textbf{Section 4.3: File-System Implementation}}\\*\hline
\textbf{4.3.1 Layout} & A volume commonly holds a boot area, superblock or volume metadata, free-space information, file metadata, directories, and data blocks. Exact placement is a design choice.\\\hline
\textbf{4.3.2 File allocation} & Contiguous allocation gives fast sequential and random access but suffers external fragmentation and growth problems. Linked allocation grows easily but has poor random access; tables or i-nodes centralize block pointers.\\\hline
\textbf{4.3.3 Directories on disk} & A directory maps names to metadata identifiers or stores attributes directly. Linear search is simple; hashing or indexed structures speed large-directory lookup.\\\hline
\textbf{4.3.4 Shared files} & Hard links make several directory entries refer to one file identity; symbolic links store another path. Reference counts, ownership, and deletion semantics must remain consistent.\\\hline
\textbf{4.3.5 Log-structured systems} & Changes accumulate in memory and are written as large sequential segments. A cleaner reclaims fragmented live data, trading background work for efficient write patterns.\\\hline
\textbf{4.3.6 Journaling} & Before modifying main structures, a file system records intended metadata or data changes in a journal. Recovery replays committed work and discards incomplete transactions.\\\hline
\textbf{4.3.7 Virtual file systems} & A virtual file-system layer exposes common operations while dispatching through per-file-system objects and methods, allowing different local and remote formats to share one namespace.\\\hline
\rowcolor{ice}\multicolumn{2}{l}{\hypertarget{sec-4-4}{}\textbf{Section 4.4: File-System Management and Optimization}}\\*\hline
\textbf{4.4.1 Disk-space management} & Block size balances internal fragmentation against metadata and transfer efficiency. Free blocks can be tracked by bitmaps, linked lists, grouping, or counting contiguous runs.\\\hline
\textbf{4.4.2 Backups} & Physical backups copy blocks; logical backups traverse files and directories. Full and incremental strategies trade restore complexity against backup time and space; consistency requires a stable view.\\\hline
\textbf{4.4.3 Consistency} & A checker compares directory references, metadata, allocation maps, and link counts after a crash. It repairs contradictions conservatively but cannot reconstruct every lost intention.\\\hline
\textbf{4.4.4 Performance} & Buffer caches, read-ahead, delayed writes, block clustering, and metadata placement reduce device operations. Policies must protect durability and avoid cache pollution.\\\hline
\textbf{4.4.5 Defragmentation} & Rearranging file blocks can improve locality and create larger free extents, but consumes I/O and must preserve a consistent mapping if interrupted.\\\hline
\rowcolor{ice}\multicolumn{2}{l}{\hypertarget{sec-4-5}{}\textbf{Section 4.5: Example File Systems}}\\*\hline
\textbf{4.5.1 MS-DOS} & \textcolor{legacy}{\textbf{Historical or legacy context:}} The FAT design stores cluster chains in a file-allocation table and directory entries with compact metadata. It is simple and widely compatible but scales poorly and provides limited protection.\\\hline
\textbf{4.5.2 UNIX V7} & \textcolor{legacy}{\textbf{Historical or legacy context:}} Directory entries map names to i-node numbers. Each i-node stores metadata and direct plus indirect block pointers, supporting small files efficiently and large files through additional indirection.\\\hline
\textbf{4.5.3 CD-ROM} & \textcolor{legacy}{\textbf{Historical or legacy context:}} Read-only optical formats emphasize portable directory records, fixed sectors, and compatibility. Extensions add longer names and richer UNIX-style metadata.\\\hline
\rowcolor{ice}\multicolumn{2}{l}{\hypertarget{sec-4-6}{}\textbf{Section 4.6: Research on File Systems}}\\*\hline
\textbf{Research themes} & Research addresses crash consistency, copy-on-write structures, checksums, deduplication, flash-aware placement, distributed naming, scalable metadata, and persistent-memory semantics.\\\hline
\rowcolor{ice}\multicolumn{2}{l}{\hypertarget{sec-4-7}{}\textbf{Section 4.7: Summary}}\\*\hline
\textbf{External and internal views} & Applications see files, directories, paths, and operations; implementers manage block allocation, metadata, naming, free space, recovery, caching, and device locality.\\\hline
\end{longtable}
\normalsize

\section*{Chapter Synthesis}
\begin{studybox}{Chapter Summary}
The file system turns storage devices into a persistent hierarchical namespace. Its external contract defines files, directories, attributes, and operations; its internal design maps names to metadata and content blocks while preserving consistency. Allocation, logging, journaling, caching, backup, and placement choices jointly determine performance and recoverability.
\end{studybox}

\begin{studybox}{Key Relationships}
\begin{itemize}
  \item Directory entries bind human-readable names to persistent file identities and metadata.
  \item Allocation metadata and file block pointers must agree for both correctness and free-space accounting.
  \item Caching and delayed writes improve speed but increase the amount of state that recovery mechanisms must reconcile.
  \item The virtual file-system layer separates a stable API from diverse local and remote implementations.
\end{itemize}
\end{studybox}

\begin{studybox}{Design Tradeoffs}
\begin{itemize}
  \item Contiguous placement maximizes locality but makes growth and fragmentation difficult.
  \item Large blocks reduce metadata and I/O overhead but waste more space in small files.
  \item Delayed writes increase performance but enlarge the window of data loss after failure.
  \item Journaling accelerates recovery but adds write traffic and ordering requirements.
\end{itemize}
\end{studybox}

\begin{studybox}{Common Confusions}
\begin{itemize}
  \item A file name belongs to a directory entry; the underlying file identity can have several names through hard links.
  \item A symbolic link contains a path, whereas a hard link refers directly to the same file identity.
  \item Crash consistency protects structural invariants; it does not automatically guarantee that every recent application write is durable.
\end{itemize}
\end{studybox}

\section*{Key Terms}
\small
\begin{longtable}{C N}
\rowcolor{navy}\color{white}\textbf{Term} & \color{white}\textbf{Concise definition}\\
\endfirsthead
\rowcolor{navy}\color{white}\textbf{Term} & \color{white}\textbf{Concise definition (continued)}\\
\endhead
\textbf{File descriptor} & Per-process handle used for operations on an open file.\\\hline
\textbf{Directory} & Namespace object mapping component names to files or subdirectories.\\\hline
\textbf{Absolute path} & Path resolved from the namespace root.\\\hline
\textbf{Relative path} & Path resolved from a process's working directory.\\\hline
\textbf{i-node} & Persistent metadata object containing attributes and data-block references.\\\hline
\textbf{FAT} & Table that records linked cluster chains and free clusters.\\\hline
\textbf{Hard link} & Additional directory entry referring to the same file identity.\\\hline
\textbf{Symbolic link} & Special file containing a target path.\\\hline
\textbf{Superblock} & Volume-level metadata describing a file system.\\\hline
\textbf{Journal} & Ordered log used to recover committed file-system changes.\\\hline
\textbf{Log-structured file system} & Design that writes accumulated changes in sequential segments.\\\hline
\textbf{VFS} & Common dispatch layer for multiple file-system implementations.\\\hline
\textbf{Read-ahead} & Prefetch of blocks predicted to be needed soon.\\\hline
\textbf{Incremental backup} & Backup containing changes since a prior backup baseline.\\\hline
\textbf{Internal fragmentation} & Unused space inside allocated blocks.\\\hline
\end{longtable}
\normalsize

\enlargethispage{2\baselineskip}
\begin{studybox}{Active-Recall Review}
\small
\begin{enumerate}
  \item How do sequential and random file access differ?
  \item What state does \texttt{open} establish that \texttt{read} and \texttt{write} later use?
  \item How is a relative path resolved?
  \item Which allocation method best supports direct access, and at what cost?
  \item How do direct and indirect i-node pointers scale file size?
  \item What failure problem does a journal solve?
  \item Why can delayed writeback improve speed yet weaken durability?
  \item How do hard and symbolic links differ after the original name is removed?
  \item What invariants can a file-system consistency checker compare?
\end{enumerate}
\end{studybox}

\par\medskip
{\color{navy}\bfseries Next-Chapter Connections}\par\smallskip
Chapter 5 moves below the file-system abstraction to controllers, drivers, interrupts, DMA, storage hardware, user interfaces, and power-aware I/O policy.
\normalsize
\end{document}
